% Kompilacja (dwa przebiegi, żeby powstał spis treści i zakładki PDF):
%   cd thesis
%   pdflatex -interaction=nonstopmode main.tex
%   pdflatex -interaction=nonstopmode main.tex
\documentclass[12pt,a4paper,oneside]{report}

\usepackage[utf8]{inputenc}
\usepackage[T1]{fontenc}
\usepackage[polish]{babel}
\usepackage{lmodern}
\usepackage{microtype}
\usepackage{geometry}
\usepackage{setspace}
\usepackage{booktabs}
\usepackage{array}
\usepackage{amsmath}
\usepackage[nottoc]{tocbibind}

\geometry{
  a4paper,
  top=2.5cm,
  bottom=2.5cm,
  left=3cm,
  right=2.5cm
}

\onehalfspacing
\setcounter{tocdepth}{3}
\setcounter{secnumdepth}{3}

\addto\captionspolish{%
  \renewcommand{\appendixname}{Załącznik}%
}

\DeclareRobustCommand{\api}[1]{\texttt{#1}}

\usepackage[unicode,hidelinks]{hyperref}
\pdfstringdefDisableCommands{%
  \def\texttt#1{#1}%
  \def\api#1{#1}%
}

% Ścieżki działają przy kompilacji z katalogu thesis/ oraz z katalogu repozytorium.
\IfFileExists{chapters/01_wstep.tex}{%
  \newcommand{\thesisinput}[1]{\input{#1}}%
}{%
  \newcommand{\thesisinput}[1]{\input{thesis/#1}}%
}

\hypersetup{
  pdftitle={Implementacja modularnego frameworka głębokiego uczenia w C++ z akceleracją GPU},
  pdfauthor={Tomasz Okoń},
  pdflang={pl}
}

\newcommand{\tytulpracy}{Implementacja modularnego frameworka głębokiego uczenia w~C++ z~akceleracją GPU na potrzeby klasyfikacji i~detekcji obiektów YOLO}
\newcommand{\uczelnia}{Politechnika Warszawska}
\newcommand{\wydzial}{Wydział Elektroniki i Technik Informacyjnych}
\newcommand{\kierunek}{Informatyka, specjalizacja sztuczna inteligencja}
\newcommand{\autorpracy}{Tomasz Okoń}
\newcommand{\promotorpracy}{prof. dr hab. inż. Robert Marek Nowak}

\begin{document}

\begin{titlepage}
  \centering
  {\large\uczelnia\par}
  \vspace{0.35cm}
  {\large\wydzial\par}
  \vspace{0.25cm}
  {\large\kierunek\par}
  \vspace{2.2cm}
  {\Large Praca inżynierska\par}
  \vspace{1.8cm}
  {\LARGE\bfseries \tytulpracy\par}
  \vspace{2cm}
  \begin{tabular}{@{}ll@{}}
    Autor: & \autorpracy \\[0.7em]
    Promotor: & \promotorpracy \\
  \end{tabular}
  \vfill
  {\large 2026\par}
\end{titlepage}

\tableofcontents
\clearpage

\thesisinput{front/wykaz_skrotow}

\thesisinput{chapters/01_wstep}
\thesisinput{chapters/02_podstawy_teoretyczne}
\thesisinput{chapters/03_architektura_i_implementacja}
\thesisinput{chapters/04_metodologia_i_benchmarki}
\thesisinput{chapters/05_wnioski}

\thesisinput{back/bibliografia}

\listoftables

\thesisinput{back/zalaczniki}

\end{document}
